\section{Introduction}

The last decade has seen astounding progress in the development of
artificial neural networks, under the ``deep learning'' rebranding
\cite{LeCun:etal:2015}.  In computer vision, we can now
automatically recognize thousands of objects in natural images
\cite{Russakovsky:etal:2014,Krizhevsky:etal:2017}. In the domain of
natural language, deep networks led to great progress in applications
ranging from machine translation to document understanding
\cite{Vaswani:etal:2017,Edunov:etal:2018,Devlin:etal:2019,Radford:etal:2019}. %
Deep networks that combine vision and language can generate image
captions and answer complex questions about scenes with high accuracy
\cite{Anderson:etal:2018,Zhou:etal:2020}.

These successes are attained by networks that are passively exposed to massive
amounts of text and/or images, and learn to rely on the
statistical regularities they extracted from their training data. %
% These computational models are akin to children that must learn their
% native language just by reading a
% mountain of books.
The interactive, functional aspects of language and intelligence \cite<e.g.,>{Wittgenstein:1953,Austin:1962,Searle:1969,Clark:1996,Allwood:1976,Pickering:Garrod:2004,Linell:2009,Ginzburg:Poesio:2016}
are completely ignored. By definition, an \emph{agent} cannot
learn to (inter-)\emph{act} just by being passively exposed to
lots of data (even when these data are records of
interactions). An exclusive focus on passive statistical learning has practical consequences. Despite much
exciting work in the area, deep-learning-based chatbots and dialogue
systems \cite{Serban:etal:2016,Gao:etal:2019} are still extremely limited in their
capabilities, missing on the dynamic, interactive nature of conversation
\cite{Bernardi:etal:2015}. And AI agents able to fully cooperate with humans
are still a science-fiction
dream \cite{Mikolov:etal:2016}.

The aim of developing devices capable of genuine linguistic interaction has revived interest in
studying the ``languages'' actively
developed by communities of artificial agents that must communicate in
order to succeed in their environment. Earlier
work in this mold \cite<e.g.,>{Cangelosi:Parisi:2002,Christiansen:Kirby:2003,Steels:2003b,Wagner:etal:2003,Steels:2012,Hurford:2014} %is this a word? yup!
explored very specific questions through carefully designed experimental
simulations that involved simple, largely hand-crafted agents. For example, \citeA{Batali:1998}  performed simulations in which a sender agent, given an input binary vector representing the meaning of a simple phrase (e.g., \textit{you smile}), encodes it as a sequence of characters. These are transmitted to a receiver agent, who needs to decode the original meaning (Fig.~\ref{fig:environments}(a)). The question under investigation was whether agent messages would mirror the grammatical structure encoded in the simple input phrase, as in natural language, and this turned out indeed to be the case.

Today,
generic ``deep agents'' (Fig.~\ref{fig:agents}), built out of standard
components such as convolutional and recurrent
  networks \cite{LeCun:etal:1998,Elman:1990,Hochreiter:Schmidhuber:1997}, with little or no task-specific tweaking,
are being used in simulations that go beyond what was conceivable just
a few years ago: dealing with complex scenarios involving thousands of
possible referents, that are presented in perceptually realistic
formats; engaging in self-paced multi-turn interactions; producing
long, language-like utterances (Fig.~\ref{fig:environments},
Fig.~\ref{fig:referential-game}).

% Figure Legends (250 words per legend)
% - Should always have a short, explanatory title (as well as legend).
% - Legends must explain the figure fully without reference to the text.
% - If applicable, the original source of previously published material must be acknowledged in the Figure legend.
% - Figures must be cited in the main text (box figures must be cited in the box). 

\begin{figure}[p]
  \centering
\includegraphics[width=14cm]{figures/environments}
\caption{\textbf{Examples of games and environments for emergent communication.}
(a) Emergent communication work in the pre-deep-learning era typically used symbolic data as input: \citeA{Batali:1998} presents a study where recurrent neural network agents communicate in a referential game using sequences of discrete symbols. Similar work with deep networks often uses realistic pictures as input, see Fig.~\ref{fig:referential-game} for an example. (b) More complex scenarios with deep agents: \citeA{Cao:etal:2018} study self-interested agents engaging in a multi-turn negotiation game. (c) Richer, dynamic environments: \citeA{Jaques:etal:2019} study five embodied self-interested agents engaging in multi-turn interactions while navigating in a 2D visual environment. (d) Scaling up to fully realistic scenarios: in the experiment of \citeA{Das:etal:2019}, embodied cooperative agents solve navigation challenges in a 3D environment. Images from \citeA{Jaques:etal:2019} and \citeA{Das:etal:2019} reproduced by permission.}\label{fig:environments}
\end{figure}


% Figure Legends (250 words per legend)
% - Should always have a short, explanatory title (as well as legend).
% - Legends must explain the figure fully without reference to the text.
% - If applicable, the original source of previously published material must be acknowledged in the Figure legend.
% - Figures must be cited in the main text (box figures must be cited in the box). 

\begin{figure}[p]
  \centering
  \includegraphics[width=14cm]{figures/architecture}
  \caption{\textbf{Typical neural network components of a deep agent.}
    (a) A \emph{visual processing} module (typically a convolutional
    network) converting pictures into internal distributed
    representations.  (b) A \textit{generation} component consisting
    of a recurrent neural network that produces a symbol sequence
    (in this case, $AXZ$).  (c) An
    \textit{understanding} module, that takes as input a sequence of
    units (in this case, the symbols produced by the generation
    component) and produces an internal distributed representation.  A typical
    \textit{sender} agent will first transform images into distributed
    representations with (a) and then use (b) to produce a message.
    A \textit{receiver} agent will also use (a) to transform images to
    representations, and then (c) to process the message from the
    sender in order to make a decision about the output action. In
    both cases, further layers are interspersed with the
    various components to further aid the agents' ``reasoning''
    process (e.g., the receiver might use them to combine visual and verbal information).
    \label{fig:agents}}
\end{figure}


% Figure Legends (250 words per legend)
% - Should always have a short, explanatory title (as well as legend).
% - Legends must explain the figure fully without reference to the text.
% - If applicable, the original source of previously published material must be acknowledged in the Figure legend.
% - Figures must be cited in the main text (box figures must be cited in the box). 

\begin{figure}[p]
  \centering
\includegraphics[width=10cm]{figures/referential-game}
  \caption{\textbf{The referential game of  \citeA{Lazaridou:etal:2017}.} In a referential game, successful communication is the very purpose of the game (as opposed to scenarios in which communication can help players to achieve an independent goal, such as obtaining a valuable object). Referential games have a long history in linguistics, philosophy and game theory \cite{Lewis:1969,Skyrms:2010}. In the game illustrated here, the sender network receives in input two natural images, depicting instances of two distinct categories out of about 500 (here: a dog and a car), with one of the images marked as target (here, the car). The sender processes the images with a convolutional network module and it emits one symbol (sampled from a fixed alphabet), that is given as input to the receiver network, together with the two images (in random order). If the receiver ``points'' to the correct location of the target in the image array (as it does in the figure), both agents are rewarded. The networks are trained by letting them play the game many times, and adjusting their weights based on the reward signal. No supervision is provided about the symbols to be used for communication, so that they are completely free to adapt the emergent protocol to their strategies and biases.}
  \label{fig:referential-game}
\end{figure}



In this survey, we review this recent literature on language emergence
in deep agent communities. After introducing some representative examples, we focus specifically on two lines of
investigation that are currently prominent in the field. First, the very complexity and richness of deep agents and their
environments implies that they often will succeed at communicating while using very opaque codes. This has led
to extensive analytical work on decoding the emergent protocol, in an attempt to understand its generality and its similarities to human language (if any), and to identify possible degenerate cases.

Sceond, we review studies that focus on how to make emergent language more powerful and
useful from an AI perspective. This involves, on the one hand, exploring
how a self-induced communication protocol might benefit deep networks
endowed with advanced perceptual and navigational capabilities. On the
other, researchers are studying how to let agents evolve  more
human-like languages, with the aim of establishing effective human-machine
communication.


\section{Language Emergence in Deep Agent Communities}
\label{sec:new-simulations}

%input{previous-simulations}

%% Figure Legends (250 words per legend)
% - Should always have a short, explanatory title (as well as legend).
% - Legends must explain the figure fully without reference to the text.
% - If applicable, the original source of previously published material must be acknowledged in the Figure legend.
% - Figures must be cited in the main text (box figures must be cited in the box). 

\begin{figure}[p]
  \centering
  \includegraphics[width=14cm]{figures/architecture}
  \caption{\textbf{Typical neural network components of a deep agent.}
    (a) A \emph{visual processing} module (typically a convolutional
    network) converting pictures into internal distributed
    representations.  (b) A \textit{generation} component consisting
    of a recurrent neural network that produces a symbol sequence
    (in this case, $AXZ$).  (c) An
    \textit{understanding} module, that takes as input a sequence of
    units (in this case, the symbols produced by the generation
    component) and produces an internal distributed representation.  A typical
    \textit{sender} agent will first transform images into distributed
    representations with (a) and then use (b) to produce a message.
    A \textit{receiver} agent will also use (a) to transform images to
    representations, and then (c) to process the message from the
    sender in order to make a decision about the output action. In
    both cases, further layers are interspersed with the
    various components to further aid the agents' ``reasoning''
    process (e.g., the receiver might use them to combine visual and verbal information).
    \label{fig:agents}}
\end{figure}


%% Figure Legends (250 words per legend)
% - Should always have a short, explanatory title (as well as legend).
% - Legends must explain the figure fully without reference to the text.
% - If applicable, the original source of previously published material must be acknowledged in the Figure legend.
% - Figures must be cited in the main text (box figures must be cited in the box). 

\begin{figure}[p]
  \centering
\includegraphics[width=10cm]{figures/referential-game}
  \caption{\textbf{The referential game of  \citeA{Lazaridou:etal:2017}.} In a referential game, successful communication is the very purpose of the game (as opposed to scenarios in which communication can help players to achieve an independent goal, such as obtaining a valuable object). Referential games have a long history in linguistics, philosophy and game theory \cite{Lewis:1969,Skyrms:2010}. In the game illustrated here, the sender network receives in input two natural images, depicting instances of two distinct categories out of about 500 (here: a dog and a car), with one of the images marked as target (here, the car). The sender processes the images with a convolutional network module and it emits one symbol (sampled from a fixed alphabet), that is given as input to the receiver network, together with the two images (in random order). If the receiver ``points'' to the correct location of the target in the image array (as it does in the figure), both agents are rewarded. The networks are trained by letting them play the game many times, and adjusting their weights based on the reward signal. No supervision is provided about the symbols to be used for communication, so that they are completely free to adapt the emergent protocol to their strategies and biases.}
  \label{fig:referential-game}
\end{figure}



%\subsection{Language Emergence Simulations and the Advent of Deep Learning}
%\label{sec:realistic-simulations}

% The focus that generative linguistics has put on language as an abstract
% structural system  (e.g., \cite{Chomsky:1965,Pinker:1994,Sag:etal:2003}) has
% long been complemented by research traditions emphasizing how language is also
% a \emph{tool} for interactively getting things done (e.g.,
% \cite{Wittgenstein:1953,Austin:1962,Searle:1969,Clark:1996,Allwood:1976,Pickering:Garrod:2004,Linell:2009,Ginzburg:Poesio:2016}).  Within the domain of computational cognitive science, this alternative view
% triggered a fruitful research line
% (\cite{Cangelosi:Parisi:2002,Christiansen:Kirby:2003,Steels:2003b,Wagner:etal:2003,Steels:2012,Hurford:2014})
% designing artificial communicating agents and testing under which conditions core properties of natural language would emerge as the result of communicative pressures.


% While this line of research addressed many interesting  questions, technology and data limitations restricted it to experimental simulations involving agents that were largely hand-crafted to address very specific issues in an \emph{ad-hoc} manner.% , making thus these simulations (and their conclusions) somewhat disconnected from the richness and complexity of natural language. 

% The last decade has seen astounding progress in the development of
% artificial neural networks, under their ``deep learning'' rebranding
% (\cite{LeCun:etal:2015}).  The most impressive successes in the
% language domain came from networks that, after being exposed to
% massive amounts of text, are able to tackle challenges ranging from
% machine translation to document understanding
% (\cite{Vaswani:etal:2017,Edunov:etal:2018,Devlin:etal:2019,Radford:etal:2019}). In
% computer vision, we can now automatically recognize thousands of objects in natural
% images (\cite{Russakovsky:etal:2014,Krizhevsky:etal:2017}). Deep
% networks combining vision and language can generate image captions
% and answer complex questions about scenes
% (\cite{Anderson:etal:2018,Zhou:etal:2020}).

% Spurred by these innovations, a new wave of language emergence research
% uses generic ``deep agents'' (Fig.~\ref{fig:agents}), built out of
% standard deep-learning components, such as convolutional and
% recurrent networks (\cite{LeCun:etal:1998,Elman:1990,Hochreiter:Schmidhuber:1997,Cho:etal:2014}),  in
% simulations that go beyond what was conceivable just a few years ago:
% dealing with complex scenarios involving hundreds of possible
% referents, that are presented in perceptually realistic formats;
% engaging in self-paced multi-turn interactions; producing long,
% language-like utterances (Fig.~\ref{fig:environments},
% Fig.~\ref{fig:referential-game}).

% % Figure Legends (250 words per legend)
% - Should always have a short, explanatory title (as well as legend).
% - Legends must explain the figure fully without reference to the text.
% - If applicable, the original source of previously published material must be acknowledged in the Figure legend.
% - Figures must be cited in the main text (box figures must be cited in the box). 

\begin{figure}[p]
  \centering
  \includegraphics[width=14cm]{figures/architecture}
  \caption{\textbf{Typical neural network components of a deep agent.}
    (a) A \emph{visual processing} module (typically a convolutional
    network) converting pictures into internal distributed
    representations.  (b) A \textit{generation} component consisting
    of a recurrent neural network that produces a symbol sequence
    (in this case, $AXZ$).  (c) An
    \textit{understanding} module, that takes as input a sequence of
    units (in this case, the symbols produced by the generation
    component) and produces an internal distributed representation.  A typical
    \textit{sender} agent will first transform images into distributed
    representations with (a) and then use (b) to produce a message.
    A \textit{receiver} agent will also use (a) to transform images to
    representations, and then (c) to process the message from the
    sender in order to make a decision about the output action. In
    both cases, further layers are interspersed with the
    various components to further aid the agents' ``reasoning''
    process (e.g., the receiver might use them to combine visual and verbal information).
    \label{fig:agents}}
\end{figure}


% % Figure Legends (250 words per legend)
% - Should always have a short, explanatory title (as well as legend).
% - Legends must explain the figure fully without reference to the text.
% - If applicable, the original source of previously published material must be acknowledged in the Figure legend.
% - Figures must be cited in the main text (box figures must be cited in the box). 

\begin{figure}[p]
  \centering
\includegraphics[width=10cm]{figures/referential-game}
  \caption{\textbf{The referential game of  \citeA{Lazaridou:etal:2017}.} In a referential game, successful communication is the very purpose of the game (as opposed to scenarios in which communication can help players to achieve an independent goal, such as obtaining a valuable object). Referential games have a long history in linguistics, philosophy and game theory \cite{Lewis:1969,Skyrms:2010}. In the game illustrated here, the sender network receives in input two natural images, depicting instances of two distinct categories out of about 500 (here: a dog and a car), with one of the images marked as target (here, the car). The sender processes the images with a convolutional network module and it emits one symbol (sampled from a fixed alphabet), that is given as input to the receiver network, together with the two images (in random order). If the receiver ``points'' to the correct location of the target in the image array (as it does in the figure), both agents are rewarded. The networks are trained by letting them play the game many times, and adjusting their weights based on the reward signal. No supervision is provided about the symbols to be used for communication, so that they are completely free to adapt the emergent protocol to their strategies and biases.}
  \label{fig:referential-game}
\end{figure}



% \subsection{Representative Experiments}
%\label{sec:representative-experiments}

In the representative simulations we will describe in this section, the agents are presented with a
 task, and each agent has a cost or reward function to optimize.
Agents have perfectly aligned incentives, i.e., they
 share their reward. Communication comes into play as a means to
achieve their goal. Learning generally takes place through %
% ; since the agents are optimizing a shared reward function,
% cooperation can lead to higher rewards and communication can facilitate
% exchange of information. Their communication behaviour is thus constantly
% adjusted to optimize their reward, and this learning takes place through
reinforcement learning, a set of techniques to train systems in scenarios in which the main teaching signal is \emph{reward} for succeeding or failing at a task, possibly requiring multiple actions in a potentially changing environment \cite{Sutton:Barto:1998,Mnih:etal:2015,Silver:etal:2016}. This setup offers more flexibility than standard \emph{supervised} learning (where the learning signal derives from direct comparison of the system output with the ground-truth), but it is also more challenging. Communication is emergent in the sense that, at the beginning of a simulation, the
 symbols the agents emit have no \emph{ex-ante} semantics nor pre-specified usage rules. Meaning and syntax emerge through game play. % 

\subsection{Continuous and Discrete Communication}
\label{sec:discrete-continuous}
 % box communication here
 Communication can be of two types:
  \textit{continuous}, in which agents communicate via a continuous
  vector, and \textit{discrete}, in which agents
  communicate by means of single symbols or sequences of symbols.
An example of the continuous case is the influential DIAL system of \citeA{Foerster:etal:2016}. The agents are given a continuous communication channel, making it easy to back-propagate learning signals through the whole system. A continuous vector connecting two agent networks can equivalently be seen as another activation layer in a larger
architecture encompassing the two networks, and it effectively gives each agent access to the internal states of the other network. Therefore, continuous communication turns the multi-agent system into a single large network.  
A ``vanilla'' model of discrete communication commonly used in the language emergence literature and for multi-agent coordination problems is RIAL \cite{Foerster:etal:2016}. In RIAL, communication happens through discrete symbols, thus making it impossible for agents to transmit rich error information via continuous back-propagation through each other. The only learning signal received by each agent is task reward. As such, unlike in the continuous case, and similarly to what happens in human communities, each agent treats the other(s) as part of its environment, with no access to their internal states. It is thus exactly the presence of the discrete bottleneck that makes simulations genuinely ``multi-agent''.  Moreover, communication via discrete symbols provides the symbolic scaffolding for interfacing the agents' emergent code to natural language,  which is universally discrete \cite{Hockett:1960}.
Tuning the weights of a neural network with an error signal that is back-propagated through a discrete bottleneck is a challenging technical problem. Adopting methods from reinforcement learning, agent training can take place using the REINFORCE update rule \cite{Williams:1992,Lazaridou:etal:2017}, which intuitively increases the weight for actions that resulted in a positive  reward  (proportional to their probability), and decreases them otherwise. Alternatively, discrete representations can be approximated by continuous ones during the training  phase \cite{Havrylov:Titov:2017,Jang:etal:2017,Maddison:etal:2017}.

\subsection{Representative Studies}

Since we initially focus on the comparison of
emergent codes with natural language,
we briefly review here work that considers the discrete case, coming back
to some examples of continuous communication in Section
\ref{sec:agent-coordination} below.

%% Figure Legends (250 words per legend)
% - Should always have a short, explanatory title (as well as legend).
% - Legends must explain the figure fully without reference to the text.
% - If applicable, the original source of previously published material must be acknowledged in the Figure legend.
% - Figures must be cited in the main text (box figures must be cited in the box). 

\begin{figure}[p]
  \centering
\includegraphics[width=10cm]{figures/referential-game}
  \caption{\textbf{The referential game of  \citeA{Lazaridou:etal:2017}.} In a referential game, successful communication is the very purpose of the game (as opposed to scenarios in which communication can help players to achieve an independent goal, such as obtaining a valuable object). Referential games have a long history in linguistics, philosophy and game theory \cite{Lewis:1969,Skyrms:2010}. In the game illustrated here, the sender network receives in input two natural images, depicting instances of two distinct categories out of about 500 (here: a dog and a car), with one of the images marked as target (here, the car). The sender processes the images with a convolutional network module and it emits one symbol (sampled from a fixed alphabet), that is given as input to the receiver network, together with the two images (in random order). If the receiver ``points'' to the correct location of the target in the image array (as it does in the figure), both agents are rewarded. The networks are trained by letting them play the game many times, and adjusting their weights based on the reward signal. No supervision is provided about the symbols to be used for communication, so that they are completely free to adapt the emergent protocol to their strategies and biases.}
  \label{fig:referential-game}
\end{figure}



% We focus specifically on work that assumes that communication takes
% place through a \emph{discrete} channel, that is, by means of symbols
% or sequences of symbols, as this is a universal property of human
% language (\cite{Hockett:1960}). Moreover, it is exactly the presence
%On of the first modern simulations in this topic was presented by 
%Seizing the opportunity afforded by advances in computer vision
%(\cite{Russakovsky:etal:2014,Krizhevsky:etal:2017}) and
%\textbf{reinforcement learning}
%(\cite{Mnih:etal:2015,Silver:etal:2016}),
% The early 2010s saw enormous progress in computer vision, driven by
% deep convolutional networks. Systems went from error-prone recognition
% of just a few categories to a near-human ability to classify thousands
% of objects in noisy real-life pictures
% (\cite{Russakovsky:etal:2014,Krizhevsky:etal:2017}). In parallel,
% successful application of reinforcement learning techniques to deep
% networks \cite{Mnih:etal:2015,Silver:etal:2016} made it possible to
% use them in dynamic scenarios involving discrete actions (such as
% emitting linguistic symbols), and relying on task success as the sole
% training signal. Seizing
% the opportunity afforded by these advances,
One of the first studies of language emergence in deep networks was presented
by \citeA{Lazaridou:etal:2017}, who used the referential game schematically
illustrated in Fig.~\ref{fig:referential-game}. %
%
% \textbf{AL: 'and to propagate
%   this signal through discrete bottlenecks' -> This seems a bit weird
%   here as we had discrete actions, we just didn't have continuous
%   representations/actions. It's the NN as function approximations that
%   made the revolution? MB: my point was about advances in general, but
%   even us, we used REINFORCE, right? It's true that that was possible
%   since 1992\ldots} Lazaridou and colleagues designed a modern variant
% of the classic signaling/referential games long employed in the
% linguistic, philosophical and game-theoretical literatures
% (\cite{Lewis:1969,Skyrms:2010}). A \emph{sender} agent (built out of
% various deep network modules) is exposed to a target and a distractor
% object, both represented by non-curated pictures of instances of about
% 500 distinct object categories (e.g., dogs and cars). After processing
% the images with a convolutional neural network module, the sender
% produces a single-symbol message. A \emph{receiver} agent (with a
% similar architecture to the sender) gets as input the sender message
% as well as the two images in random order, and has to point to the
% intended target. If it points to the right target, both agents receive
% positive reward, which is the only training signal.
This paper first
showed that agents can develop an effective communication protocol to talk about realistic
images by relying on game success as sole training signal. Still, evidence that the agents were developing human-like words referring to generic concepts such as ``dog'' or ``animal'' was mixed,
a point will return to in the next section.

While \citeA{Lazaridou:etal:2017} constrained messages to consist of one
symbol, \citeA{Havrylov:Titov:2017} allowed the sender to emit strings
of symbols of variable length
\cite<see also>{Lazaridou:etal:2018}. The resulting emergent language developed
a prefix-based hierarchical scheme to encode meaning into
multiple-symbol sequences. For example, the ``word'' for pizza was
\emph{5261 2250 5211}, where \emph{5261} refers to food, \emph{2250}
to baked food, and \emph{5211} to pizzas. %However, unlike real words, symbol
%meanings tended to change with their position in the message. 

\citeA{Evtimova:etal:2018} went one step further, considering multiple-turn
interactions \cite<see also>{Jorge:etal:2016}.
% Human linguistic
% interaction is obviously not limited to single-message episodes, so
% another move towards more realistic simulations involves allowing
% agents to take multiple turns. A more flexible scenario was explored
% by Evtimova et al.~\cite{Evtimova:etal:2018}.
In their game, one agent
must pick the definition of an animal from a list of dictionary
entries, when a natural image of the target animal is presented to
the other agent. %The messages can include multiple symbols,
% transmitted in an order-insensitive way.
The agents can exchange
multiple messages. The conversation ends when
the agent tasked with guessing the definition makes its
final guess.  Several natural
properties of conversations emerged in this setup. For example, the
agents tend to exchange more turns in more difficult game
episodes. % Moreover, message information content increases with turns
% for the image describing agent, whereas it decreases for the
% definition guessing one. This stands to reason: as the game
% progresses, the definition guesser can rely on accrued information,
% and ask more specific questions. On the other hand, these questions
% will require more nuanced answers on behalf of the image describer.



% Jorge et
% al.~\cite{Jorge:etal:2016} first studied this in a highly constrained
% setup where the agents where assigned a fixed number of single-symbol
% turns.


% Going beyond classic referential games, Cao et
% al.~\cite{Cao:etal:2018} (partially inspired by
% \cite{Lewis:etal:2017}) studied whether agents can learn to use
% multiple-turn conversations to achieve their goals in a negotiation
% scenario. In each episode of their game, two agents A and B are
% exposed to a (symbolically represented) set of items. Each of them is
% also provided with (random and hidden) utilities for each item type
% (e.g., in a specific episode, peppers might be very valuable and
% cherries useless for agent A). At each step, agents emit a
% multiple-symbol message, as well as a (non-verbally conveyed) proposal
% on how to split the goods. Either agent can terminate the episode at
% any time by accepting the proposal that the other agent issued at the
% previous step. If the agents do not reach an agreement within 10
% turns, neither gets any reward. The agents learn to play the game at
% least somewhat successfully in all setups, but the crucial factor
% determining whether they will meaningfully use the linguistic channel
% (as opposed to only rely on the information directly conveyed by the
% non-verbal proposals) is whether they are rewarded ``prosocially''
% (each agent receives the cumulative reward of both agents under the
% agreed split) or egotistically. Only prosocial agents learn to use
% language, sharing in particular information about their
% utilities. When one agent is made to play against a community of other
% agents, the language channel, again, is meaningfully used only when
% there are enough prosocial agents in the community. However, each pair
% of agents develops its own idiolect, with no convergence to a common
% language.

A further step towards realistic conversational scenarios, beyond
referential games, is taken by \citeA{Bouchacourt:Baroni:2019} \cite<see also>{Cao:etal:2018}. In this study, one agent is assigned a
fruit, the other two tools, and their task is to decide which of the
two tools is best for the current fruit. The utility of each tool with
respect to each fruit is derived from a corpus of human judgments,
resulting in skewed affordance statistics (e.g., a knife is generally
more useful than a spoon). The setup is fully symmetric, with either
agent randomly assigned either role in each episode, and both agents
being able to start and end the conversation. The agents learn to use
messages meaningfully, accumulating more reward than what they could
get by relying on general object affordances. However, despite the
symmetric setup, they develop different idiolects for the different
roles they take, that is, the same agents use different codes to
communicate the same meanings, depending on who is in charge of describing the fruit, and who the tools. A similar behaviour  was also observed by \citeA{Cao:etal:2018}.
% one role or the other, start the conversation or go second, and decide
% whether to end the episode. Going beyond classic referential games,
% one agent is assigned a fruit, the other two tools (symbolically
% represented), and their task is to cooperatively decide which tool is
% best for the current fruit. The utility of each tool with respect to
% each fruit is derived from a corpus of human judgments, resulting in
% skewed affordance statistics (e.g., a knife is generally more useful
% than a spoon). Bouchacourt and Baroni show that the agents learn to
% use messages meaningfully, being able in this way to accumulate more
% reward than what they could get by relying on general object
% affordances. Moreover, unlike in earlier studies such as
% \cite{Evtimova:etal:2018}, where artificially ``lobotomizing'' one of
% the agents was a pre-condition for communication to help, in this
% experiment the agents rely on language even when they are both endowed
% with full memory capabilities. On the other hand, despite the strongly
% symmetric setup, the agents develop different idiolects that they use
% when taking different roles (a result also reported by
% \cite{Cao:etal:2018}).

Under which conditions will agents converge to a shared language is
one of the topics addressed by \citeA{Graesser:etal:2019}, who used deep agent communities as a
modeling tool for contact linguistics
\cite{MyersScotton:2002}. They report that,
if two agents will develop different idiolects, it suffices for
the community to include a third agent for a shared code to
emerge. One of their most interesting findings is that, when agent
communities of similar size are put in contact, the agents develop a
mixed code that is simpler than either original language, akin
to the development of pidgins and creoles in mixed-language
communities \cite{Bakker:etal:2011}.
% Two
% agents are exposed to different fragments of a synthetic image The
% agents communicate through order-insensitive multi-symbol messages for
% a fixed number of turns. At the end, they must both select the correct
% caption for the whole image (e.g., ``there is a green triangle'') out
% of a set of alternatives. If only two agents are playing the game,
% they will develop different idiolects. It suffices however for three
% agents to play the game (in turns) for a shared code to
% emerge. Moreover, increasing community size does not slow down
% learning. Graesser and colleagues then study what happens when agents
% from two communities that were originally disconnected are made to
% interact. They find that all agents are soon able to communicate
% across communities, including those that were never in direct contact
% with agents from the other community. When communities of
% different sizes are made to interact, the smaller community will adopt
% the code of the larger one. However, and perhaps most excitingly from
% the point of view of contact-linguistics studies, if the communities
% are of equal size, the agents will develop a mixed code that is
% simpler than either of the original languages, akin to the development of
% pidgins and creoles in mixed-language communities.

% As we have seen, different studies took different steps towards
% greater realism. However, to the best of our knowledge, nobody has yet
% tried to combine all these advances in a single simulation that would
% feature, say, natural visual inputs, multi-turn dialogue, variable
% length messages, more than 2 agents, and a task based on real-world
% statistics. This might in part be due to the technical complications
% that such pursuit implies. However, a possibly even more important
% factor is that, already in the simpler experiments we reviewed here,
% decoding the agents' language, and understanding exactly how it is being
% used, turns out to be a major challenge. 
% Thus, many recent studies are
% actually considering even \emph{simpler} setups, and focusing on
% analysis of the emergent code. This is the topic we turn to in the
% next section.



\section{Understanding the Emergent Language}
\label{sec:understanding-emergent-language}

%\textbf{AL: I'm now wondering if we should have a 3-section piece, with understanding/measuring being its own thing after the new simulations and the agent coordination?}
With more realistic simulations, understanding
what is going on becomes more difficult. Even when we are confident
that genuine communication is taking place (Section
\ref{sec:measuring-communication}), it is difficult to decode messages
by simple inspection. We do not know if and how the messages produced
by the agents should be segmented into ``words''. We might only have
vague conjectures about what they refer to. If there are multiple
turns, we do not know which turns are mostly information exchanges,
and which, if any, absolve other pragmatic functions (e.g., asking for more information). We cannot even
trust the agents to use symbols in a consistent way across contexts
and turns \cite{Bogin:etal:2018}. The enterprise is akin to
linguistic fieldwork, except that we are dealing with an alien race,
with no guarantees that universals of human communication will
apply.% There has thus been extensive work focusing
% on decoding the agents' emergent language.

Indeed, in spite of task success, the emerging language can have
counter-intuitive properties. \citeA{Kottur:etal:2017} considered
agents playing a referential game in which they must communicate about
object attributes and values (e.g, \emph{color: blue}, \emph{shape:
  round}). The agents have difficulties converging to the intuitive
coding scheme in which distinct symbols unambiguously denote single
attributes or values (i.e., a word for \emph{color}, a word for
\emph{blue}, etc). Such code will only emerge when the set of
available symbols is greatly limited and the memory of one of the
agents is ablated, pointing to memory bottlenecks as a possible
bias to be injected into deep networks for more natural languages to
emerge \cite<see also>{Resnick:etal:2020}. \citeA{Bouchacourt:Baroni:2018} replicated
the game of \citeA{Lazaridou:etal:2017} with the surprising results
illustrated in Fig.~\ref{fig:gaussian}.
% considered a multi-turn game based on
% single-symbol exchanges. An agent sees an object, and the other is
% tasked with reporting the values of two attributes of the object
% (e.g., color and shape). If the agents are allowed to use a large set
% of symbols, they will learn to associate symbols to whole attribute
% value combinations (e.g., a single symbol meaning
% \emph{blue-square}). This strategy results in a language that cannot
% generalize to new combinations of attributes or values. More
% interestingly, even when the available symbol cardinality is reduced,
% the agents do not converge to the intuitive strategy of using them to
% denote single attributes and values. Instead, the agent that needs to
% ask about attributes saves on vocabulary usage by letting the same
% symbol refer to attribute combinations in an order-insensitive way
% (e.g., one symbol denotes both \emph{color-shape} and
% \emph{shape-color}). The other agent uses symbol sequences to refer to
% different value combinations in an arbitrary, context-sensitive way
% (e.g., symbol \emph{A} followed by \emph{B} might refer to
% \emph{blue}, but when followed by \emph{C} it might refer to
% \emph{square}). Kottur et al.~finally show that, only when the memory
% of one agent is removed, so that the latter cannot rely on context,
% the emergent language words consistently denote single attributes and
% values. This comes, however, at the cost of agent lobotomy. Bogin et
% al.~\cite{Bogin:etal:2018} also found that, with standard training
% methods and a sufficiently complex environment, agents will not
% develop a consistent language where symbols have the same referent
% independently of context. 

% Figure Legends (250 words per legend)
% - Should always have a short, explanatory title (as well as legend).
% - Legends must explain the figure fully without reference to the text.
% - If applicable, the original source of previously published material must be acknowledged in the Figure legend.
% - Figures must be cited in the main text (box figures must be cited in the box). 

\begin{figure}[p]
  \centering
\includegraphics[width=7cm]{figures/gaussian}
\caption{\textbf{Training and test inputs in the referential game of
    \cite{Bouchacourt:Baroni:2018}.} Two agents were trained to play
  the  game of \citeA{Lazaridou:etal:2017} (see
  Fig.~\ref{fig:referential-game}). During training, the agents were
  exposed to the same data as in the original study, that is, pairs of
  pictures of instances of about 500 distinct objects (top row). At
  test time, however, the agents were made to play the game with blobs
  of Gaussian noise (bottom row). They were able to
  communicate about them nearly as well as about the training
  pictures. This shows that the language emerging in this game does
  not involve ``words'' referring to generic concepts, but rather
  \emph{ad-hoc} signals, probably carrying comparative information
  about shallow visual properties of the images. Bottom row reproduced
  from \citeA{Bouchacourt:Baroni:2018} by permission.}
  \label{fig:gaussian}
\end{figure}


Essentially, agents will develop a code that is sufficient to solve
the task at hand, and hoping that such code will possess further
desirable characteristics is wishful thinking. Consider the task in
the original game of \citeA{Lazaridou:etal:2017}. The agents must
discriminate pairs of pictures depicting instances of 500
categories. The agents could achieve this by developing human-like
names for the categories, but a low-level strategy relying on, say,
comparing average pixel intensity in patches of the two images might
require as few symbols as 2. In this respect, the
agents' language is, paradoxically, ``too human'', in the sense that
it evolved to minimize effort, while remaining adequate for the task
at hand \cite{Gibson:etal:2019}. % ; compare to the tendency of humans in
% restricted communication setups to converge to \emph{ad-hoc} naming
% conventions, or ``conceptual pacts'',
% \cite{Brennan:Clark:1996}). 
Indeed, \citeA{Kharitonov:etal:2020} showed
that the way deep agent emergent languages partition their
meaning space displays the same tendency towards complexity minimization
that is pervasive in human language.

%Still on the topic of effort minimization,
\citeA{Chaabouni:etal:2019} studied whether agent language
% In part because of the aforementioned difficulties, other studies
% focus on arguably even more basic properties of natural language than
% compositionality. For example, it is well known that human languages
% and other animal communication systems
exhibit an inverse correlation between word frequency and word length,
so that the signals that need to be used more often are also the
shortest, as universally found in natural languages
\cite{Zipf:1949,Strauss:etal:2007,FerrerICancho:etal:2013}. They discovered that deep agents trained with a referential game where inputs
have a skewed distribution similar to natural language actually
develop a significantly \emph{anti-efficient} code, in which the
most frequent inputs are associated to the longest messages. The
effect is explained by the lack of an articulatory effort
minimization bias in networks, that are thus only subject to a
``perceptual'' pressure favoring longer messages, as they are easier to discriminate.
% would also develop a language exhibiting this
% property. Not only this was not the case: the agents developed a
% positively \emph{anti-efficient} code, in which the most frequent
% inputs were associated to the longest messages. They showed that this
% surprising pattern emerges because listening agents have an innate
% preference for longer messages, as they are more discriminative. The
% same might be true of human and animal perceivers. However, in
% biological agents this is counterbalanced by the speaker preference
% for conciseness, as message articulation is costly. Since artificial
% networks are not subject to comparable metabolic constraints, there is
% no similar contrasting pressure, and  codes that only favour the listener can
% easily come about. Paradoxically, by providing a counterexample to its universal
% emergence, this finding supports the view that the inverse
% frequency/length correlation attested in humans and animal
% communication systems is indeed due to efficiency factors. If it was a
% trivial effect of random string assignment processes, as sometimes
% conjectured (\cite{Miller:1957,delPradoMartin:2013}), then it should
% arise in any successful communication system dealing with a skewed
% input referent distribution.

% As another example, Chaabouni et al.~\cite{Chaabouni:etal:2019} studied whether deep agents have innate
% word order biases. The agents were trained to use simple languages
%  displaying or lacking certain properties. Relative ease of learning by
% individual agents, as well as language drift in cross-generation
% transmission, were used to probe how natural the relevant
% characteristics were for the agents. Chaabouni and colleagues found
% deep agents to have priors in line with some universal preferences of
% human language, such as a tendency to reduce word order variation, and
% avoidance of long-distance dependencies. Intriguingly, agents, like
% humans, prefer symbols referring to a sequence of events to be encoded
% iconically, but in \emph{reverse} order (that is, saying something
% like \emph{turn left, go straight} to mean ``first go straight, then
% turn left''). Finally, in another example of how artificial agents are
% not subject to the same least-effort pressures that shape human
% language, they found that the agents did not have a tendency to shed
% redundant cues of the same information (such as double role marking
% by means of word order and overt suffixes). All in all, the results of this study further confirm the mixed picture about the relation between deep agent codes (and the underlying biases that lead to their emergence) and natural language.  

% box communication
\subsection{Measuring the Degree of Effective Communication}
\label{sec:measuring-communication}
  As simulations move beyond referential games (where task success
  trivially depends on establishing a communication code), to complex
  environments where communication plays an auxiliary function (e.g.,
  \citeA{Das:etal:2019}; Fig.~\ref{fig:environments}(d)), the first
   question to ask when analyzing an emergent language is whether
  it is actually been used in any meaningful way by the agents.  As
  clearly discussed by \citeA{Lowe:etal:2019}, just ablating the
  language channel and showing a drop in task success does not prove
  much, as the extra capacity afforded by the channel architecture
  might have helped the agents' learning process without being used to
  establish communication. The same paper proposes a classification of
  measures to detect the presence of genuine
  communication. \textit{Positive signaling} captures the extent to
  which information about the sender states, observations and actions
  are expressed in its signals. \textit{Positive listening} captures
  the extent to which a signal impacts the receiver's states and
  behaviour. Examples of positive signaling include \textit{context
    independence}~\cite{Bogin:etal:2018} and \textit{speaker
    consistency}~\cite{Jaques:etal:2019}. The former measures the
  degree of alignment between messages and task-related concepts,
  whereas the latter measures, through mutual information, the
  alignment between an agent messages and its actions. Positive
  signaling gives no guarantee of communication, since the receiver could
  be ignoring the sender messages, no matter how informative they might
  be. An example of positive listening is the \textit{instantaneous
    coordination} measure of~\citeA{Jaques:etal:2019}, which uses
  mutual information to quantify correlation between 
  sender messages and receiver actions. Instead,~\citeA{Lowe:etal:2019}
  propose to use \emph{causal influence of communication}, a
  quantity that measures the causal relationship between sender
  messages and receiver actions.   The authors show that only a high causal influence
  of communication is both necessary and sufficient for
  positive listening, and thus communication.

  The call for caution
  is not just hypothetical. Several studies have reported how agents
  can easily converge to non-verbal or degenerate strategies, even
  when it would seem that communication is taking place. For example,
   agents might learn to exchange information simply through the
  number of turns they take before ending the game, irrespective of
  what they actually say
  \cite{Cao:etal:2018,Bouchacourt:Baroni:2019}.

\subsection{Compositionality}

% Starting from similar considerations, much work on analyzing emergent
% languages has focused on which environmental aspects and biases
% (``innate'' or injected into the agent architectures) force the
% development of a code possessing some properties of interest.

Much analytical work in the area has focused on compositionality, as the latter is seen both as a
fundamental feature of natural language whose evolutionary origins are
unclear \cite{Bickerton:2014,Townsend:etal:2018}, and as a pre-condition
for an emergent language to generalize at scale.
% There
% has been in particular widespread interest in the conditions under
% which the emergent language is compositional. This interest is
% theoretically motivated by the idea that compositionality is a
% fundamental, perhaps unique feature of human language, but its
% evolutionary origins are not clear
% (\cite{Bickerton:2014,Townsend:etal:2018}). From an applied angle,
% compositionality might be a pre-condition for a language to be able
% to generalize at scale.

% In linguistics, compositionality is traditionally defined as the
% property by which the meaning of a composite linguistic expression is
% a function of the meanings of its parts (and systematic meaning
% composition rules) (\cite{Partee:2004}). In the context of current
% language emergence work, composite meanings are typically restricted
% to simple conjunctions of attribute values (is the message denoting
% \emph{red triangles} predictable from messages denoting other
% \emph{red} things and messages denoting \emph{triangles} of other
% colors?). Even so, it turns out to be very hard to assess how
% compositional an emergent language is. First, we have not yet
% developed systematic techniques to segment messages into component
% parts (note that there is no guarantee that such parts should be
% combined by agents concatenatively). Second, a black-or-white view of compositionality, such as the
% one suggested by the standard linguists' definition, is too coarse. It
% is likely that different emergent languages will show different
% degrees of compositionality, and missing this gradient might imply missing
% important insights onto what makes a language take the first
% steps towards full compositionality.

The simplest way to probe for compositionality in an emergent protocol
is to test whether agents can use it to denote novel composite
meanings, e.g., can they refer to \emph{blue squares} on first encounter, if
they have seen other \emph{blue} and \emph{square} things during
training \cite{Choi:etal:2018}. This assumes that a
compositional encoding is necessary to generalize. However, a few recent papers have intriguingly reported %
% address these challenges is to directly test the
% ability of agents to refer to novel composite meanings as a
%  proxy of compositionality (e.g., can the agents
% successfully refer to a \emph{blue square} the first time they are
% exposed to it, if in the training phase they encountered other
% \emph{blue} and \emph{square} things?)
% (e.g., \cite{Choi:etal:2018}). However, this approach is assuming that
% something akin to the linguists' notion of compositionality is
% necessary for generalization, which is actually the sort of hypothesis
% we would like to \emph{validate} with our simulations. Indeed, one of
% the most intriguing observations coming from the field is
that emergent languages can support generalization to novel composite meanings \emph{without} conforming to even weak notions of
compositionality \cite{Lazaridou:etal:2018,Andreas:2019,Chaabouni:etal:2020}. %  an
% empirical result that we hope to be explored in much more detail in
% the future. Moreover, generalization might tell us at best that some
% form of compositionality is in place, but it won't give us any cue on
% the kind of compositional processes that a language is making use of.
\citeA{Lazaridou:etal:2018} re-introduced to the emergent language
community the \emph{topographic similarity} score from earlier work on
language emergence \cite{Brighton:2006}. Given ways to measure
distances between meanings and between forms, topographic similarity
is the correlation between all possible meaning pair distances and the
distances of the corresponding message pairs. It captures the intuition that compositionality involves a systematic
relation between form and meaning similarity. However, it does not tell us
anything about the nature of the specific compositional processes present
in a language. %
% and it does so by
% providing a continuous score. Moreover, minimum edit distance is
% relatively robust to non-concatenative changes in form, as it can
% capture ways to compose messages involving deletions and
% insertions. Conversely, exactly because of its generality, a high
% topographic similarity does not tell us much about the nature of the
% compositional processes that are present in a language.
\citeA{Andreas:2019}
proposed a method to quantify to what extent an emergent language reflects specific types of compositional structure in its
input. Unfortunately, the method only works if we have a concrete
hypothesis about the underlying composition function, that is, it can only be used to test whether a language conforms to an underlying compositional grammar if we are able to precisely specify this grammar. This limits its practical
applicability. Finally, \citeA{Chaabouni:etal:2020} recently established a link between compositionality and the notion of disentanglement in representation learning \cite{Suter:etal:2019}, and proposed to use methods to quantify disentanglement from that literature in order to measure the degree of compositionality of emergent codes.

Equipped with similar tools, various studies have uncovered different aspects
of compositionality in emergent languages. For example, \citeA{Lazaridou:etal:2018}
found that compositionality more easily emerges when objects are
represented symbolically as sets of attribute-value pairs, than when
they are more realistically represented as synthetic 3D shapes. %
%
% let two agents play a referential game where the speaker was presented
% with a target object and produced a (multi-symbol) message, and the
% receiver had to point to the target amidst distractors. They found
% that when objects were represented as symbolic attribute-value sets,
% the emergent language exhibited at least a certain degree of
% compositional structure, as reflected by topographic similarity. The
% picture was more mixed when realistic perceptual processing was added
% to the simulation, with objects being represented as 3D geometrical
% shapes in synthesized scenes, that were given to the agents as raw
% pixel maps \textbf{(add figure to better explain this?)}. On the one
% hand, some degree of compositionality was also present in at least
% some of the experiments run under this regime. However, the pattern
% was now quite brittle. In particular, reducing the number of
% distractors led to the emergence of non-compositional,
% hard-to-interpret languages, probably focusing on low-level aspects of
% the comparison between the target and the few distractor images
% left. Again, the take-home message might be that intuitive languages
% (in this case, languages characterizing geometric shapes in terms of
% discrete visual attributes) will only emerge if the task is
% challenging enough to justify the complexity of the solution (the more
% compositional language that emerged in the simulations uses more than
% 1000 distinct messages, compared to a dozen or less for the degenerate
% ones). As Lazaridou and colleagues point out, the tendency to converge
% to \emph{ad-hoc} solutions that are good-enough for the task at hand
% is not an unnatural property of deep agents, but rather an instance of
% the same efficiency pressure that leads humans to establish
% specialized naming conventions (``conceptual pacts'') in specific
% communicative settings (\cite{Brennan:Clark:1996}).
% Other studies have explored other aspects of compositionality in
% emergent languages. For example,
% Mordatch and Abbeel
\citeA{Mordatch:Abbeel:2018}  studied
the code emerging in a community of agents moving and
acting in a shared grid-world. Each agent was assigned a goal, that
could involve having another agent moving to a landmark position. An
order-insensitive concatenative language emerged, where agents would
refer to actions, their agent and targets by juxtaposing specialized
symbols (e.g., one message could be: \emph{goto blue-agent
  red-landmark}, or, equivalently, \emph{red-landmark goto
  blue-agent}). % When agents were deprived of the
% communication channel, they resorted to non-verbal strategies, such as
% pointing or pushing, to direct each other.

Despite many intriguing empirical observations, our characterization
of which architectural biases and environmental pressures favour the
emergence of compositionality (or other linguistic properties) is
still very sketchy. A strong result obtained both with humans and in
pre-deep-learning computational simulations is that generational
transmission of language favors compositionality
\cite<e.g.,>{Kirby:Hurford:2002,Kirby:etal:2014}, an observation recently
confirmed for deep agents \cite{Li:Bowling:2019,Ren:etal:2019}.
Moreover, recent results from experiments with humans demonstrate that
larger communities of speakers evolve more systematic languages
\cite{Raviv:etal:2019a,Raviv:etal:2019b}, suggesting the need to
move away from two-partner agent setups, an observation
beginning to find its way into deep agent research
\cite{Graesser:etal:2019,Tieleman:etal:2019}.

Other priors for compositionality that have been proposed and at least partially
empirically validated include input representations
\cite{Lazaridou:etal:2018}, agent and channel capacity
\cite{Kottur:etal:2017,Mordatch:Abbeel:2018,Resnick:etal:2020}, and
specific training strategies, such as letting the agents simulate
other agents' understanding of one's language
\cite{Choi:etal:2018}. We still lack, however, systematic
experiments establishing which of these conditions are necessary,
which are sufficient, and how they interact, possibly along the lines
of earlier systematizing work such as \citeA{Bratman:etal:2010}.
% Overall, though, our characterization of the emergence of compositionality in
% deep agent communication is much murkier than in the small controlled
% experiments from the previous decades. First, a number of possible
% priors favoring compositionality have found preliminary support,
% including, besides generational transmission, constraints on the %environmental pressures
% environment and input representations (\cite{Lazaridou:etal:2018}),
% agent and channel capacity
% (\cite{Kottur:etal:2017,Mordatch:Abbeel:2018}), and specific training
% strategies, such as letting the agents simulate other agents'
% understanding of one's language (\cite{Choi:etal:2018}). We still
% lack systematic experiments establishing which
% of these conditions are necessary, which are sufficient, and how they
% interact. Second, except in the tiniest experiments, we never really
% observe the emergence of a fully compositional language. Rather, emergent languages tend at best to mix messages (or message substrings) that can be
% transparently segmented into components and messages that are
% partially or completely opaque. Further analysis of this intriguing
% result is needed.


